\chapter{Graph Theory}\label{graph-theory}
\chapterstatus{complete}

\section{Introduction}\label{graph-theory:section-introduction}

A graph is the simplest mathematical model of a relation between objects: a set of \emph{vertices} and a set of \emph{edges}, each edge joining two vertices. Road and railway networks,
molecules, social networks, the moves of a game, the dependencies between the results of this book (which the website draws as a graph) are all graphs, and many questions about them,
whether one can get from anywhere to anywhere, how robustly, whether a tour visiting every place once exists, how to pair up objects, how many resources are needed to avoid conflicts,
have a precise graph-theoretic formulation. Graph theory began with Euler's solution of the problem of the bridges of K\"onigsberg in 1736 (Section~\ref{graph-theory:section-euler-hamilton}).

This chapter develops the structural theory of finite graphs:
\begin{enumerate}
\item the basic notions (subgraphs, degrees, isomorphism) and the handshake lemma, with the characterisation of degree sequences (Section~\ref{graph-theory:section-graphs});
\item walks, paths, cycles, connectedness, distance and bipartite graphs (Section~\ref{graph-theory:section-paths});
\item trees: their many equivalent definitions, spanning trees, Pr\"ufer's proof of Cayley's formula and minimum spanning trees (Section~\ref{graph-theory:section-trees});
\item Euler tours and Hamilton cycles: Euler's theorem, de Bruijn sequences, the theorems of Dirac, Ore and Bondy--Chv\'atal (Section~\ref{graph-theory:section-euler-hamilton});
\item matchings: Berge's augmenting paths, the theorems of Hall, K\"onig, Tutte and Petersen, and stable matchings (Section~\ref{graph-theory:section-matchings});
\item connectivity: blocks, ear decompositions, Menger's theorem and the max-flow min-cut theorem (Section~\ref{graph-theory:section-connectivity});
\item colourings: greedy colouring and degeneracy, Brooks' theorem, graphs of large chromatic number without triangles, the chromatic polynomial, edge colourings (K\"onig and Vizing),
and an overview of perfect graphs (Section~\ref{graph-theory:section-colourings}).
\end{enumerate}
Planar graphs, extremal and Ramsey theory and spectral graph theory have chapters of their own (Chapters~\ref{planar-graphs}, \ref{extremal-graph-theory}
and~\ref{spectral-graph-theory}); random graphs appear in Chapter~\ref{probabilistic-method}. We follow mostly \cite{diestel-graph-theory}, and also \cite{bondy-murty} and \cite{west-graph-theory};
\cite{lovasz-plummer} is the reference for matchings.

\noindent\textbf{Prerequisites.} Chapter~\ref{sets} (Sets) and Chapter~\ref{combinatorics} (Combinatorics); Section~\ref{graph-theory:section-trees} compares with
Chapter~\ref{generating-functions} (Generating Functions).

\begin{notation}\label{graph-theory:notation-basic}
As in Chapter~\ref{combinatorics}, $[n] = \{1, \ldots, n\}$, $|A|$ is the number of elements of a finite set $A$, and $\binom{A}{2}$ is the set of $2$-element subsets of $A$. A partition of a
set into two pieces is written $A \sqcup B$. Graphs are finite unless said otherwise; the conventions for writing edges, degrees and subgraphs are fixed in
Notation~\ref{graph-theory:notation-graphs}.
\end{notation}

\section{Graphs}\label{graph-theory:section-graphs}

\subsection{Definitions and first examples}\label{graph-theory:subsection-definitions}

\begin{definition}\label{graph-theory:definition-graph}
A \emph{graph} is a pair $G = (V, E)$ of a finite set $V$, the set of \emph{vertices}, and a set $E \subseteq \binom{V}{2}$ of $2$-element subsets of $V$, the \emph{edges}. We write $V(G) = V$ and
$E(G) = E$. Two vertices $u, v$ are \emph{adjacent}, or \emph{neighbours}, if $\{u, v\} \in E$; the edge $\{u, v\}$ \emph{joins} $u$ and $v$, which are its \emph{ends}, and is \emph{incident}
with them. The \emph{order} of $G$ is $|V|$, and its \emph{size} is $|E|$.
\end{definition}

\begin{notation}\label{graph-theory:notation-graphs}
We write the edge $\{u, v\}$ as $uv$ (so $uv = vu$), and $u \sim v$ when $uv \in E$. For a vertex $v$:
\begin{enumerate}
\item $N(v) = N_G(v) = \{u \in V : uv \in E\}$ is its \emph{neighbourhood}, and $d(v) = d_G(v) = |N(v)|$ is its \emph{degree};
\item $\delta(G) = \min_vd(v)$ and $\Delta(G) = \max_vd(v)$ are the \emph{minimum} and \emph{maximum degree} (for $V \neq \emptyset$), and $\bar{d}(G) = \frac{1}{|V|}\sum_vd(v)$ is the
\emph{average degree};
\item a vertex of degree $0$ is \emph{isolated}, and a vertex of degree $1$ is a \emph{leaf}; $G$ is \emph{$k$-regular} if $d(v) = k$ for every $v$, and \emph{cubic} if it is $3$-regular.
\end{enumerate}
For a set $S \subseteq V$, $N(S) = \bigcup_{v \in S}N(v)$. When several graphs are around we write $N_G$, $d_G$, and we always write $n = |V|$ and $m = |E|$ when one graph is fixed.
\end{notation}

The definition excludes \emph{loops} (an edge from a vertex to itself) and \emph{multiple edges} (several edges with the same ends): these are the \emph{simple graphs}, and they are what
``graph'' means in this chapter. The excluded objects are useful too (Definition~\ref{graph-theory:definition-multigraph}). A graph is usually drawn with points for the vertices and curves
for the edges; the drawing is not part of the graph, and the same graph has many drawings (Example~\ref{graph-theory:example-petersen}).

\begin{example}[Standard graphs]\label{graph-theory:example-standard}
\begin{enumerate}
\item The \emph{complete graph} $K_n$ has vertex set $[n]$ and all $\binom{n}{2}$ edges; it is $(n - 1)$-regular. The \emph{empty graph} $\overline{K_n}$ on $[n]$ has no edges.
\item The \emph{path} $P_n$ ($n \geq 1$) has vertices $1, \ldots, n$ and edges $\{i, i + 1\}$, $1 \leq i < n$; it has $n - 1$ edges, two leaves when $n \geq 2$, and \emph{length} $n - 1$.
\item The \emph{cycle} $C_n$ ($n \geq 3$) has vertices $1, \ldots, n$ and the edges of $P_n$ together with $\{n, 1\}$; it is $2$-regular with $n$ edges. $C_3 = K_3$ is the \emph{triangle}.
\item The \emph{complete bipartite graph} $K_{r,s}$ has vertex set $A \sqcup B$ with $|A| = r$, $|B| = s$, and the $rs$ edges $ab$, $a \in A$, $b \in B$. The \emph{star} $K_{1,s}$ has a
\emph{centre} joined to $s$ leaves.
\item The \emph{$d$-dimensional cube} $Q_d$ has vertex set $\{0, 1\}^d$, two words being adjacent if they differ in exactly one coordinate; it is $d$-regular with $2^d$ vertices and
$d2^{d-1}$ edges.
\item For a group $\Gamma$ and a subset $S \subseteq \Gamma$ with $1 \notin S$ and $S = S^{-1}$, the \emph{Cayley graph} has vertex set $\Gamma$ (assumed finite) and the edges $\{g, gs\}$,
$s \in S$; it is $|S|$-regular. The cycle $C_n$ is the Cayley graph of $\ZZ/n\ZZ$ with $S = \{\pm 1\}$, and $Q_d$ that of $(\ZZ/2\ZZ)^d$ with $S$ the standard basis.
\end{enumerate}
\end{example}

\begin{example}[The Petersen graph]\label{graph-theory:example-petersen}
The \emph{Petersen graph} $P$ has as vertices the $10$ two-element subsets of $[5]$, two being adjacent when they are disjoint. Each $2$-subset is disjoint from $\binom{3}{2} = 3$ others, so
$P$ is cubic with $15$ edges. It is usually drawn as a pentagon $\{1,2\}, \{3,4\}, \{5,1\}, \{2,3\}, \{4,5\}$ (consecutive ones disjoint), a pentagram inside it, and five spokes joining
each outer vertex $\{a, b\}$ to the inner vertex $\{c, d\}$ disjoint from it. The Petersen graph is the source of many counterexamples in graph theory: it is cubic, has no cycle of length
$3$ or $4$ (two disjoint $2$-subsets of $[5]$ have exactly one common neighbour... in fact none, see Exercise~\ref{graph-theory:exercise-petersen}), and is not Hamiltonian
(Example~\ref{graph-theory:example-petersen-not-hamiltonian}).
\end{example}

\begin{definition}\label{graph-theory:definition-isomorphism}
An \emph{isomorphism} $G \to H$ of graphs is a bijection $\phi \colon V(G) \to V(H)$ with $uv \in E(G) \iff \phi(u)\phi(v) \in E(H)$ for all $u, v$. Graphs are \emph{isomorphic},
$G \cong H$, if there is an isomorphism between them. An isomorphism $G \to G$ is an \emph{automorphism}; the automorphisms form a group $\operatorname{Aut}(G)$ under composition, a subgroup of the
symmetric group of $V(G)$. A \emph{graph property} is a property of graphs that is invariant under isomorphism.
\end{definition}

Isomorphic graphs differ only in the names of their vertices; in pictures, one usually draws an \emph{unlabelled} graph, that is, an isomorphism class. The number of vertices, the number of
edges, the multiset of degrees, connectedness and the number of triangles are graph properties; the statement ``vertex $1$ is a leaf'' is not. The number of graphs on the vertex set $[n]$ is
$2^{\binom{n}{2}}$, while the number of isomorphism classes was computed for $n = 4$ in Example~\ref{combinatorics:example-graphs-four-vertices}: there are $11$.

\begin{example}\label{graph-theory:example-automorphisms}
$\operatorname{Aut}(K_n) = \operatorname{Aut}(\overline{K_n}) = S_n$; $\operatorname{Aut}(C_n)$ is the dihedral group of order $2n$ (the rotations and reflections of a regular $n$-gon),
since an automorphism of $C_n$ is determined by the image of the vertex $1$ and of its neighbour $2$, which must be a neighbour of the image of $1$; and $\operatorname{Aut}(P)$, for the Petersen
graph, contains $S_5$, acting on $2$-subsets of $[5]$ (it is in fact equal to $S_5$, Exercise~\ref{graph-theory:exercise-petersen}).
\end{example}

\begin{counterexample}\label{graph-theory:counterexample-same-degrees}
Graphs with the same degree multiset need not be isomorphic. The disjoint union of two triangles and the cycle $C_6$ are both $2$-regular on $6$ vertices with $6$ edges, but the
second is connected and the first is not. Similarly, $K_{3,3}$ and the \emph{triangular prism} (two triangles $abc$ and $a^{\prime}b^{\prime}c^{\prime}$ joined by $aa^{\prime}$, $bb^{\prime}$, $cc^{\prime}$)
are connected cubic graphs on $6$ vertices, but the prism contains triangles and $K_{3,3}$ does not (Proposition~\ref{graph-theory:proposition-bipartite}). No simple list of numerical
invariants is known that decides isomorphism.
\end{counterexample}

\subsection{Subgraphs and operations}\label{graph-theory:subsection-subgraphs}

\begin{definition}\label{graph-theory:definition-subgraph}
A graph $H$ is a \emph{subgraph} of $G$, written $H \subseteq G$, if $V(H) \subseteq V(G)$ and $E(H) \subseteq E(G)$. It is a \emph{spanning} subgraph if $V(H) = V(G)$, and an
\emph{induced} subgraph if $E(H) = E(G) \cap \binom{V(H)}{2}$; the induced subgraph on $U \subseteq V(G)$ is written $G[U]$. We write:
\begin{enumerate}
\item $G - U = G[V \setminus U]$ for $U \subseteq V$, and $G - v = G - \{v\}$;
\item $G - F = (V, E \setminus F)$ for $F \subseteq E$, $G - e = G - \{e\}$, and $G + e = (V, E \cup \{e\})$ for $e \in \binom{V}{2}$;
\item $\overline{G} = (V, \binom{V}{2} \setminus E)$ for the \emph{complement} of $G$;
\item $G \cup H = (V(G) \cup V(H), E(G) \cup E(H))$, and $G \sqcup H$ when $V(G) \cap V(H) = \emptyset$ (the \emph{disjoint union}).
\end{enumerate}
A set of pairwise adjacent vertices is a \emph{clique}, and a set of pairwise non-adjacent vertices is an \emph{independent} (or \emph{stable}) set; $\omega(G)$ and $\alpha(G)$ are the
largest sizes of a clique and of an independent set. Cliques of $G$ are independent sets of $\overline{G}$, so $\omega(G) = \alpha(\overline{G})$.
\end{definition}

\begin{example}\label{graph-theory:example-subgraphs}
In $C_5$ the path $P_5$ obtained by deleting one edge is a spanning subgraph that is not induced; deleting a vertex gives the induced subgraph $P_4$. The complement of $C_5$ is again a
$5$-cycle (join each vertex to the two vertices at distance $2$), so $C_5$ is \emph{self-complementary}; so is $P_4$. Since $\binom{n}{2}$ must be even for a self-complementary graph on
$n$ vertices (the graph and its complement have the same number of edges), $n \equiv 0$ or $1 \pmod 4$.
\end{example}

\begin{definition}\label{graph-theory:definition-multigraph}
A \emph{multigraph} is a triple $(V, E, \iota)$ of finite sets $V$, $E$ and a map $\iota$ assigning to each $e \in E$ a set $\iota(e) \subseteq V$ with one or two elements, its ends; edges with one
end are \emph{loops}. The degree of $v$ counts the edges at $v$, each loop twice. A \emph{directed graph} (or \emph{digraph}) is a pair $(V, A)$ with $A \subseteq V \times V$ minus the diagonal;
$(u, v) \in A$ is an \emph{arc} from its \emph{tail} $u$ to its \emph{head} $v$, and $d^+(v)$, $d^-(v)$ are the numbers of arcs with tail, resp.\ head, $v$ (\emph{out-} and \emph{in-degree}).
\end{definition}

Most statements of this chapter are about graphs; when multigraphs or digraphs are needed (Euler tours, flows, edge colourings), the text says so.

\subsection{Degrees and the handshake lemma}\label{graph-theory:subsection-degrees}

\begin{proposition}[Handshake lemma]\label{graph-theory:proposition-handshake}
For every graph, $\sum_{v \in V}d(v) = 2|E|$. Consequently the number of vertices of odd degree is even, and $\delta(G) \leq \bar{d}(G) = \frac{2|E|}{|V|} \leq \Delta(G)$. The same formula holds for
multigraphs (loops counted twice), and for digraphs $\sum_vd^+(v) = \sum_vd^-(v) = |A|$.
\end{proposition}

\begin{proof}
Count the pairs $(v, e)$ with $v \in V$, $e \in E$ and $v \in e$ in two ways (double counting, Section~\ref{combinatorics:section-double-counting}): for each $v$ there are $d(v)$ of them, and
for each edge there are $2$. Then $\sum_vd(v)$ is even, and the sum of the even degrees is even, so the sum of the odd degrees is even, which forces an even number of odd terms. The inequality
is that between a minimum, a mean and a maximum. For multigraphs, a loop contributes $2$ to the degree of its end, and for digraphs each arc has one tail and one head.
\end{proof}

The name comes from the reading: at a party, the total number of hands shaken is even. The lemma is used constantly; for instance, a cubic graph has an even number of vertices, and
there is no graph with degrees $3, 3, 3, 1$ (Example~\ref{graph-theory:example-havel-hakimi}).

\begin{definition}\label{graph-theory:definition-degree-sequence}
The \emph{degree sequence} of a graph with vertices $v_1, \ldots, v_n$ is the sequence $d(v_1) \geq d(v_2) \geq \cdots \geq d(v_n)$, for an ordering of the vertices by decreasing
degree. A sequence $d_1 \geq \cdots \geq d_n$ of nonnegative integers is \emph{graphic} if it is the degree sequence of some graph.
\end{definition}

\begin{theorem}[Havel--Hakimi]\label{graph-theory:theorem-havel-hakimi}
Let $n \geq 2$ and $d_1 \geq d_2 \geq \cdots \geq d_n \geq 0$ be integers with $d_1 \leq n - 1$. The sequence $(d_1, \ldots, d_n)$ is graphic if and only if the sequence
\[
(d_2 - 1, d_3 - 1, \ldots, d_{d_1+1} - 1, d_{d_1+2}, \ldots, d_n),
\]
obtained by deleting $d_1$ and subtracting $1$ from the $d_1$ next largest terms, is graphic (after reordering).
\end{theorem}

\begin{proof}
\emph{If.} Given a graph on $v_2, \ldots, v_n$ with $d(v_i) = d_i - 1$ for $2 \leq i \leq d_1 + 1$ and $d(v_i) = d_i$ for $i > d_1 + 1$, add a vertex $v_1$ joined to $v_2, \ldots, v_{d_1+1}$.
\emph{Only if.} Among all graphs on $v_1, \ldots, v_n$ with $d(v_i) = d_i$, choose one in which the number $t$ of neighbours of $v_1$ among $S = \{v_2, \ldots, v_{d_1+1}\}$ is as large as
possible. We show that $t = d_1$, that is $N(v_1) = S$; then deleting $v_1$ gives a graph realising the second sequence. Suppose $t < d_1$. Then there are $v_i \in S$ with $v_1v_i \notin E$
and $v_j \notin S$, $j \geq 2$, with $v_1v_j \in E$. Since $i < j$, $d_i \geq d_j$; but $v_j$ has the neighbour $v_1$ which $v_i$ lacks, so $v_i$ has a neighbour $w \neq v_1$ that is not a
neighbour of $v_j$ (otherwise $N(v_i) \subseteq N(v_j) \setminus \{v_1\}$ and $d_i < d_j$). Replace the edges $v_1v_j$ and $v_iw$ by $v_1v_i$ and $v_jw$ (a \emph{$2$-switch}): the degrees are
unchanged, the new edges were not present ($v_1v_i \notin E$, $v_jw \notin E$), and $v_1$ has gained the neighbour $v_i \in S$ and lost $v_j \notin S$, so $t$ increases, a contradiction.
\end{proof}

\begin{example}\label{graph-theory:example-havel-hakimi}
Is $(4, 4, 3, 3, 2, 2)$ graphic? Apply the theorem repeatedly: $(4, 4, 3, 3, 2, 2) \to (3, 2, 2, 1, 2) = (3, 2, 2, 2, 1) \to (1, 1, 1, 1) \to (0, 1, 1) = (1, 1, 0) \to (0, 0)$, and
$(0, 0)$ is realised by two isolated vertices; reversing the steps builds a graph. On the other hand $(3, 3, 3, 1) \to (2, 2, 0) \to (1, -1)$ fails, so $(3, 3, 3, 1)$ is not graphic, although its
sum is even: the parity condition of the handshake lemma is necessary but not sufficient.
\end{example}

The proof gives more: any two graphs with the same degree sequence on the same vertex set are connected by a sequence of $2$-switches (Exercise~\ref{graph-theory:exercise-two-switches}).
A closed criterion is the theorem of Erd\H{o}s and Gallai: $d_1 \geq \cdots \geq d_n$ with even sum is graphic if and only if $\sum_{i=1}^{k}d_i \leq k(k - 1) + \sum_{i=k+1}^{n}\min(d_i, k)$ for
all $k$ (necessity is Exercise~\ref{graph-theory:exercise-erdos-gallai}; for sufficiency see \cite[Section~3.1]{west-graph-theory}).

\section{Walks, paths and cycles}\label{graph-theory:section-paths}

\subsection{Walks and connectedness}\label{graph-theory:subsection-walks}

\begin{definition}\label{graph-theory:definition-walk}
A \emph{walk} of length $k \geq 0$ in $G$ is a sequence $W = (v_0, v_1, \ldots, v_k)$ of vertices with $v_{i-1}v_i \in E$ for $1 \leq i \leq k$; it goes \emph{from} $v_0$ \emph{to} $v_k$ (a
\emph{$v_0$--$v_k$ walk}) and is \emph{closed} if $v_0 = v_k$. It is a \emph{trail} if its $k$ edges $v_{i-1}v_i$ are distinct, and a \emph{path} if its $k + 1$ vertices are distinct. A \emph{cycle}
of length $k \geq 3$ is a closed walk $(v_0, \ldots, v_k = v_0)$ whose vertices $v_0, \ldots, v_{k-1}$ are distinct. We identify a path or a cycle with the subgraph it traverses (isomorphic to
$P_{k+1}$, resp.\ $C_k$), and say that $G$ \emph{contains} a path or cycle of length $k$ if it has such a subgraph.
\end{definition}

\begin{lemma}\label{graph-theory:lemma-walk-path}
Every $u$--$v$ walk contains a $u$--$v$ path: there is a $u$--$v$ path whose vertices and edges are among those of the walk. Every closed walk of odd length contains a cycle of odd length.
\end{lemma}

\begin{proof}
Induction on the length $k$ of the walk $(v_0, \ldots, v_k)$. If its vertices are distinct, it is a path. Otherwise $v_i = v_j$ for some $i < j$, and $(v_0, \ldots, v_i, v_{j+1}, \ldots, v_k)$ is a
shorter $u$--$v$ walk with vertices and edges among the old ones. For the second claim, let $W$ be a closed walk of odd length $k$, minimal among the closed odd walks with edges in $W$. If
$v_i = v_j$ for some $0 \leq i < j < k$, then $W$ splits into the two closed walks $(v_i, \ldots, v_j)$ and $(v_j, \ldots, v_k = v_0, \ldots, v_i)$ of lengths $j - i$ and $k - (j - i)$, one of
which is odd and shorter. So $v_0, \ldots, v_{k-1}$ are distinct, and as $k$ is odd and a closed walk of length $1$ is impossible (no loops), $k \geq 3$: $W$ is an odd cycle.
\end{proof}

\begin{counterexample}\label{graph-theory:counterexample-even-walk}
The second statement fails for even lengths: the closed walk $(u, v, u)$ of length $2$ contains no cycle at all, and in a tree (Section~\ref{graph-theory:section-trees}) there are closed walks of
every even length and no cycles.
\end{counterexample}

\begin{definition}\label{graph-theory:definition-connected}
A graph is \emph{connected} if it is nonempty and any two vertices are joined by a path. The relation ``$u$ and $v$ are joined by a walk'' is an equivalence relation on $V$ (reflexive by walks of
length $0$, symmetric by reversing, transitive by concatenating); its classes are the vertex sets of the \emph{components} of $G$, the induced subgraphs on the classes. We write $c(G)$ for the
number of components.
\end{definition}

By Lemma~\ref{graph-theory:lemma-walk-path} it makes no difference whether walks or paths are used. The components of $G$ are exactly its maximal connected subgraphs, and $G$ is the disjoint
union of its components: an edge $uv$ joins two vertices of the same class.

\begin{proposition}\label{graph-theory:proposition-connected-edges}
\begin{enumerate}
\item $G$ is connected if and only if for every partition $V = A \sqcup B$ into nonempty sets there is an edge with one end in $A$ and one in $B$.
\item Deleting an edge $e$ increases the number of components by at most $1$: $c(G) \leq c(G - e) \leq c(G) + 1$. Consequently a graph with $n$ vertices and $m$ edges has at least $n - m$
components, and a connected graph on $n$ vertices has at least $n - 1$ edges.
\end{enumerate}
\end{proposition}

\begin{proof}
(1) If $G$ is connected and $a \in A$, $b \in B$, a path from $a$ to $b$ has a first vertex in $B$, and the edge entering it crosses. Conversely, if $G$ is not connected, take $A$ a component and $B$
its complement: no edge crosses.
(2) The components of $G - e$, $e = uv$, other than those of $u$ and $v$ are components of $G$; the component of $G$ containing $e$ is the union of the components of $u$ and $v$ in $G - e$, since
every vertex of it is joined to $u$ by a path in $G$, which uses $e$ at most once. Starting from $\overline{K_n}$ ($n$ components) and adding the $m$ edges one at a time, the number of components
drops by at most $1$ each time.
\end{proof}

\begin{definition}\label{graph-theory:definition-bridge}
An edge $e$ is a \emph{bridge} if $c(G - e) > c(G)$, and a vertex $v$ is a \emph{cut vertex} if $c(G - v) > c(G)$.
\end{definition}

\begin{proposition}\label{graph-theory:proposition-bridge}
An edge is a bridge if and only if it lies on no cycle.
\end{proposition}

\begin{proof}
Let $e = uv$. If $e$ lies on a cycle, the rest of the cycle is a $u$--$v$ path in $G - e$, and any walk in $G$ using $e$ can be rerouted along it, so $c(G - e) = c(G)$. Conversely, if $e$ is not a
bridge then $u$ and $v$ are in the same component of $G - e$ (by the proof of Proposition~\ref{graph-theory:proposition-connected-edges}(2)), and a $u$--$v$ path in $G - e$ together with $e$
is a cycle through $e$ (the path has length at least $2$, since its only possible edge of length $1$ would be $e$).
\end{proof}

\subsection{Distance, diameter and girth}\label{graph-theory:subsection-distance}

\begin{definition}\label{graph-theory:definition-distance}
The \emph{distance} $d_G(u, v)$ is the length of a shortest $u$--$v$ path, or $\infty$ if there is none. The \emph{eccentricity} of $v$ is $\max_ud(u, v)$, the \emph{diameter}
$\operatorname{diam}(G)$ is the largest eccentricity and the \emph{radius} $\operatorname{rad}(G)$ the smallest. The \emph{girth} $g(G)$ is the length of a shortest cycle ($\infty$ if $G$ has
no cycle), and the \emph{circumference} is the length of a longest cycle.
\end{definition}

On a connected graph, $d_G$ is a metric on $V$ with integer values: $d(u, v) = 0$ only for $u = v$, it is symmetric, and the triangle inequality holds by concatenating paths
(Lemma~\ref{graph-theory:lemma-walk-path}). Also $\operatorname{rad}(G) \leq \operatorname{diam}(G) \leq 2\operatorname{rad}(G)$, the second by the triangle inequality through a vertex of
minimal eccentricity.

\begin{proposition}\label{graph-theory:proposition-bfs}
Let $G$ be connected and $r \in V$. Let $L_i = \{v : d(r, v) = i\}$ (the \emph{levels}, computed by breadth-first search). Every edge joins two vertices of the same level or of consecutive
levels, and every $v \in L_i$, $i \geq 1$, has a neighbour in $L_{i-1}$.
\end{proposition}

\begin{proof}
If $uv \in E$, then $|d(r, u) - d(r, v)| \leq 1$ by the triangle inequality. If $d(r, v) = i \geq 1$, the vertex before $v$ on a shortest $r$--$v$ path has distance $i - 1$ from $r$, since a shorter
path to it would give a shorter path to $v$.
\end{proof}

\begin{proposition}\label{graph-theory:proposition-min-degree}
Let $G$ be a graph with $\delta(G) = \delta \geq 2$. Then $G$ contains a path of length $\delta$ and a cycle of length at least $\delta + 1$. If moreover $G$ has girth $g$, then, with
$r = \lfloor(g - 1)/2\rfloor$,
\[
|V| \geq 1 + \delta\sum_{i=0}^{r-1}(\delta - 1)^i \qquad \text{(the Moore bound)}.
\]
\end{proposition}

\begin{proof}
\emph{Paths and cycles.} Let $P = (v_0, \ldots, v_k)$ be a longest path. All neighbours of $v_k$ lie on $P$, since otherwise $P$ could be extended. So $k \geq d(v_k) \geq \delta$. If $v_i$ is the
neighbour of $v_k$ with smallest index, then $i \leq k - \delta$, and $(v_i, \ldots, v_k, v_i)$ is a cycle of length $k - i + 1 \geq \delta + 1$.
\emph{Moore bound.} Fix a vertex $r_0$ and let $L_i$ be the levels from $r_0$. For $i < r$, each vertex $v \in L_i$ ($i \geq 1$) has exactly one neighbour in $L_{i-1}$ and none in $L_i$:
otherwise two shortest paths from $r_0$, or a shortest path and an edge inside a level, would close a cycle of length at most $2i + 1 \leq 2r - 1 < g$. Indeed, two different neighbours
$u, u^{\prime} \in L_{i-1}$ of $v$, with shortest paths to $r_0$, give a closed walk of length $2i$ through $v$ that, followed from $v$ until the two paths meet, contains a cycle of length at most
$2i \leq 2r - 2 < g$. So the vertices of $L_i$, $1 \leq i \leq r - 1$, have at least $\delta - 1$ neighbours in $L_{i+1}$, and these are distinct for distinct vertices of $L_i$ (a vertex of
$L_{i+1}$ with two neighbours in $L_i$ would close a cycle of length at most $2(i + 1) \leq 2r < g$, as $2r \leq g - 1$). Hence $|L_1| \geq \delta$ and $|L_{i+1}| \geq (\delta - 1)|L_i|$
for $i \leq r - 1$, and $|V| \geq \sum_{i=0}^{r}|L_i| \geq 1 + \delta\sum_{i=0}^{r-1}(\delta - 1)^i$.
\end{proof}

For $\delta = 3$ and $g = 5$ ($r = 2$) the bound is $1 + 3 + 3 \cdot 2 = 10$, attained by the Petersen graph: it is the unique smallest cubic graph of girth $5$. Graphs attaining the Moore bound with
odd girth are rare: for $g = 5$ they exist only for $\delta = 2, 3, 7$ and possibly $57$ (the Hoffman--Singleton theorem, proved with the eigenvalues of the adjacency matrix,
Chapter~\ref{spectral-graph-theory}).

\begin{counterexample}\label{graph-theory:counterexample-min-degree}
The bounds of Proposition~\ref{graph-theory:proposition-min-degree} are sharp and need $\delta \geq 2$: $K_{\delta+1}$ has minimum degree $\delta$, its longest paths have length $\delta$ and its
longest cycles length $\delta + 1$. A graph with $\delta = 1$ need not contain any cycle (a path), and large average degree does not by itself bound the \emph{minimum} degree: a star $K_{1,s}$ has
average degree close to $2$ but $\delta = 1$. (Every graph of average degree $d$ has a subgraph of minimum degree greater than $d/2$, Exercise~\ref{graph-theory:exercise-dense-subgraph}.)
\end{counterexample}

\subsection{Bipartite graphs}\label{graph-theory:subsection-bipartite}

\begin{definition}\label{graph-theory:definition-bipartite}
A graph is \emph{bipartite} if $V = A \sqcup B$ with every edge having one end in $A$ and one in $B$; $(A, B)$ is a \emph{bipartition}, and $A$, $B$ are its \emph{sides} (independent sets).
More generally $G$ is \emph{$k$-partite} if $V$ is a union of $k$ independent sets.
\end{definition}

\begin{proposition}[K\"onig]\label{graph-theory:proposition-bipartite}
A graph is bipartite if and only if it contains no cycle of odd length.
\end{proposition}

\begin{proof}
A closed walk in a bipartite graph alternates between the sides, so it has even length; in particular there are no odd cycles. Conversely, suppose there is no odd cycle. It suffices to
bipartition each component, so let $G$ be connected, fix $r$, and put $A = \{v : d(r, v) \text{ even}\}$, $B = \{v : d(r, v) \text{ odd}\}$. If an edge $uv$ had both ends in the same set, then
$d(r, u) = d(r, v)$ by Proposition~\ref{graph-theory:proposition-bfs}, and shortest paths from $r$ to $u$ and to $v$ with the edge $uv$ form a closed walk of odd length $2d(r, u) + 1$, which
contains an odd cycle by Lemma~\ref{graph-theory:lemma-walk-path}.
\end{proof}

The proof is an algorithm: breadth-first search either produces a bipartition or finds an edge inside a level, hence an odd cycle. Trees, even cycles, cubes $Q_d$ (colour a word by the parity of
its number of $1$s) and grid graphs are bipartite; odd cycles, $K_n$ for $n \geq 3$ and the Petersen graph (which contains $5$-cycles) are not.

\begin{proposition}\label{graph-theory:proposition-bipartite-subgraph}
Every graph $G$ has a bipartite spanning subgraph with at least $|E|/2$ edges.
\end{proposition}

\begin{proof}
Among all partitions $V = A \sqcup B$ choose one maximising the number of edges between $A$ and $B$. Every vertex $v$ then has at least half its neighbours on the other side: otherwise moving $v$ to
the other side would increase the number of crossing edges by (neighbours on its own side) $-$ (neighbours on the other side) $> 0$. Summing over $v$, the crossing edges, counted twice, are at least
half of $\sum_vd(v) = 2|E|$.
\end{proof}
